\scanpage{29}
\mspage{25}
% archive sheet number: 28

\uline{Digression} 1\quad Somme directe de magasins.

Si $(\unsure{M_i|\mathcal M_i})_{i\in I}$ est une famille de magasins, on définit le magasin somme $M$ ainsi
\[
M=\coprod_i M_i
\tag*{a)}
\]
b) $\trianglelefteq$ sur $M$ est la relation d'ordre somme

c) $\overset{\circ}{\ll}$ sur $M$ est la relation d'ordre somme

d) Si $x,y\in M$, on a
\[
x\mathrel{|\circ|}y\Longleftrightarrow
\left\{\begin{aligned}
&x\in M_i,y\in M_j,\ i\neq j\\
&\text{\uline{ou}}\\
&x,y\in M_i,\quad x\mathrel{|\circ|}y\\
&\text{au sens de }M_i
\end{aligned}\right.
\]

Vérification des axiomes tautologique.

Si pour chaque $i\in I$, on a choisi $\unsure{\mathcal F_i|\mathscr F_i|F_i}\in\unsure{P(M_i)|\mathcal P(M_i)}$, atelier associé à $M_i$, on définit
\[
\begin{aligned}
&\mathcal F\subset P(M)\\
&\rotatebox{90}{$=$}\\
&\{A\subset M\mid\forall i\in M_i,\quad A\cap M_i\in\mathcal F_i\}
\end{aligned}
\]
On voit que c'est un atelier associé à $M$, appelé \uline{atelier somme} des ateliers $(M_i,\trianglelefteq,\overset{\circ}{\ll},|\circ|,\mathcal F_i)$. On note que
\[
\mathcal F\simeq\prod_{i\in I}\mathcal F_i
\]
isomorphisme respectant les ordres $\unsure{\leq|\ll}$ $\subseteq,\ll$.
